\documentclass[a4paper,10pt]{article}

%%%%%%%%%%%% PREÁMBULO %%%%%%%%%%%%%%%%%%%%%

% Paquetes

\usepackage[utf8]{inputenc}
\usepackage[spanish,es-tabla]{babel}
\usepackage[T1]{fontenc}

\usepackage{listings}
\usepackage[pdftex]{graphicx} % graficos
\usepackage{multirow, array} % para las tablas
\usepackage{float} % para usar [H]
\usepackage{booktabs}
\graphicspath{ {imagen/} }%En que carpeta estan las imagenes
\usepackage{subcaption}
\usepackage{caption}


\DeclareGraphicsExtensions{.pdf,.jpeg,.png}
\usepackage{epstopdf}

% Comandos


\renewcommand{\lstlistingname}{Código}
\renewcommand{\lstlistlistingname}{Índice de fragmentos de código fuente}

% Opciones

\title{Descripción DataSet*}
\author{Viviana Mercado}

%%%%%%%%%%%%%%%%%%%%%%%%%%%%%%%%%%%%%%%%%%%%%%

\begin{document}
\maketitle

\begin{abstract}
Este documento es una pequeña guía de la conformación conjunto DataSet de artículos periodísticos 
\end{abstract}

\tableofcontents
\listoftables 	%Indices de Tablas
\listoffigures 	%Indices de Figuras


\section{Datos}

En este trabajo, se centrará en un un conjunto de documentos periodísticos argentinos extraídos de blog de noticias, redacciones, libros. 

Se ha generado un corpus de 200 documentos de 10 periodistas (5 oficialistas y 5 opositores) representándose los documentos con un enfoque de bolsa de términos (bag of terms) variándose en cada caso el tipo de término utilizado (n-gramas de palabras, n-gramas de caracteres, términos “lematizados”, términos “truncados”, etc.). 
En cada caso se han utilizado distintos enfoques para el pesado de los términos como la ponderación binaría, ponderación tf-idf, etc. 
En particular, han representado los documentos utilizando todas las características del sistema Language Inquiry Word Count (LIWC).

Después del proceso de recopilación, y debido a la heterogénea naturaleza en las codificaciones de conjuntos de caracteres de textos recopilados (ASCII, UTF-8, ANSI e ISO-8859-1), el formato de los documentos fue unificado convirtiéndolos a la misma codificación.
Se presenta el corpus del texto en español en una versión balanceada estandarizando el texto en UTF-8.

\subsection{Recopilación de datos}
Las versiones equilibradas poseen las siguientes características: 

\begin{itemize}
	\item Los Textos fueron escritos por hispanohablantes de Argentina 
	\item Los textos “hablan” hacen hincapié en aspectos políticos.
	\item Todos los documentos incluyen texto “formal” es decir, sin (o con una cantidad mínima de) errores tipográficos, contracciones, jergas (slangs), hipervínculos, etiquetas, figuras y emoticones.
	\item Hay un mínimo de 18 y un máximo de 20 documentos periodísticos tomados desde blog, textos periodísticos, libros (archivo) por autor.
	\item Para la detección de la orientación se dividió en 2 clases (oficialistas/opositores)
	\item Igual número de periodistas masculinos y femeninos en cada categoría.
 
\end{itemize}
Después de realizar el proceso de recolección, todos los documentos fueron etiquetados manualmente teniendo en cuenta la información proporcionada por los datos revelados dentro del texto.

	\begin{table}[htbp]
	\centering
	%\caption{Distribución de documentos de la versión balanceada }
	\begin{tabular}{lrrrr}
		\multicolumn{1}{c}{Autores (Docs) } & \multicolumn{1}{c}{Diarios} & \multicolumn{1}{c}{Blogs } & \multicolumn{1}{c}{Libros} & \multicolumn{1}{c}{Total} \\
		\midrule
		Oficialistas & 50    & 46    & 2     & 98 \\
		Opositores & 60    & 36    & 2     & 98 \\
		\bottomrule
	\end{tabular}%
	\label{tab:tabla1}%
	\caption{Distribución de documentos de la versión balanceada }
	\end{table}

\subsection{Análisis}
 Se presenta Tabla 2 donde muestra las características principales de ambas categorías (oficialistas /opositores) del corpus. Esto incluye el número total de palabras, oraciones, número de palabras únicas encontradas en todos Documentos (vocabulario) y el tamaño de cada versión (kilobytes)
 
 \begin{table}[htbp]
 	\centering
 	 	\begin{tabular}{lrrrr}
 		\multicolumn{1}{c}{Autores (Docs)} & \multicolumn{1}{c}{Palabras} & \multicolumn{1}{c}{Oraciones} & \multicolumn{1}{c}{Vocabulario} & \multicolumn{1}{c}{Tamaño } \\
 		\midrule
 		Oficialistas & 113.016   &   2.197,58    &   7.804,05    & 1.133 \\
 		Opositores &  172.234     &    2.353,72   & 7.883,69      & 1.586 \\
 		\bottomrule
 	\end{tabular}%
 	\label{tab:tabla2}%
\caption{Características del Corpus}
 \end{table}%

 En la Tabla 3, Figura 1 y 2 describen las características de palabras y oraciones en el corpus en la categoría oficialistas. 
 
 \begin{table}[htbp]
 	\centering
 	\begin{tabular}{lcccc}
	\cmidrule{2-5}  
	  \multicolumn{1}{r}{} & \multicolumn{1}{l}{Promedio} & \multicolumn{1}{l}{Desv.Est.} & Máx.   & Mín. \\
	\midrule
		Palabras & 1757,49 & 4866,35 & 37069,00 & 126,00 \\
		Oraciones &  24,02& 10,41 & 54,50 & 8,09 \\
		\bottomrule
	\end{tabular}%
	\caption{\label{tab:tabla3} Distribución de Palabras y Oraciones de autores oficialistas}%
\end{table}

%%%%%
% COLOCAR FIGURA 1
%%%%%

\begin{figure}%[H]
	\begin{minipage}[b]{0.5\linewidth}
	\centering
	\includegraphics[width=\linewidth]{figura1pal_of.eps}
	\caption{Palabras autores oficialistas}
	\label{fig:figura1}
	\end{minipage}
	\hspace{0.2cm}
\begin{minipage}[b]{0.5\linewidth}
	\centering
	\includegraphics[width=\linewidth]{figura2ora_of.eps}
	\caption{Oraciones autores oficialistas}
	\label{fig:figura1}
\end{minipage}
\hspace{0.5cm}

\end{figure}
\hspace{0.5cm}
%%%%
 En cambio, Tabla 4, Figuras 3 y 4 describen las características de palabras y oraciones en el corpus en la categoría opositores.
%%%%

\begin{table}[htbp]
	\centering
	\begin{tabular}{lcccc}
	\cmidrule{2-5}  
	\multicolumn{1}{r}{} & \multicolumn{1}{l}{PROMEDIO} & \multicolumn{1}{l}{DESV.EST.} & MÁX   & MÍN \\
		\midrule
	Palabras & 1153,22 & 679,36 & 3171,00 & 126,00 \\
	Oraciones & 22,42 & 5,53 & 40,62 & 12,78 \\
	\bottomrule
	\end{tabular}%
	\caption{\label{tab:tabla4}Distribución de Palabras y Oraciones de autores opositores}%
\end{table}%
\hspace{0.1cm}
%%%%%
%COLOCAR FIGURA 2
%%%%%%
\vspace{-0.5cm}
\begin{figure}%[H]
		\begin{minipage}[b]{0.5\linewidth}
		\centering
		\includegraphics[width=\linewidth]{figura3pal_op.eps}
		\caption{Palabras autores opositores}
	\label{fig:figura3}
		\end{minipage}
	\hspace{0.2 cm}
		\begin{minipage}[b]{0.5\linewidth}
		\centering
		\includegraphics[width=\linewidth]{figura4ora_op.eps}
		\caption{Oraciones autores opositores}
		\label{fig:figura4}
	\end{minipage}
	\hspace{0.5cm}
\end{figure}

\newpage
\subsection{Clases}

De la combinación de los documento, se pueden derivar dos clases. Estos se identificaron como: oficialistas y opositor. En tablas 3 y 4, se muestra estadísticas las versiones para cada clase. Incluyen, para cada combinación (categoría), el número de de documentos, número total de palabras, número promedio de palabras y oraciones respectivamente, vocabulario (número de palabras diferentes) y tamaño (en KiloBytes). De la Tabla 3 podemos observar que en la versión equilibrada, los autores oficialistas usan más palabras, más oraciones y tienen un vocabulario más rico que las autores opositores.
 
\end{document}